% Generated by scripts/assemble.py; corrections belong in source/corrections.json.
\examhead{147}{湖南师范大学 2024 年数学分析试卷}{139}
% p140-b22; source PDF page 140
\par\Needspace{5\baselineskip}
\blockanchor{p140-b22}
\majorquestion{一}填空题. 每题 6 分, 共 60 分.\par
% p140-b23; source PDF page 140
\blockanchor{p140-b23}
\question{1}极限  $\lim_{n\to\infty}\frac{1}{n}\left(\sin\left(\frac{n+1}{2n}\pi\right)+\sin\left(\frac{n+2}{2n}\pi\right)+\cdots+\sin\pi\right)=$  \_\_\_\_.\par
% p141-b00; source PDF page 141
\blockanchor{p141-b00}
\question{2}极限  $\lim_{x\to0}\frac{\int_{0}^{x}t\sin t\,dt}{\sin x-\tan x}=$  \_\_\_\_.\par
% p141-b01; source PDF page 141
\blockanchor{p141-b01}
\question{3}设 $f(x) = \frac{1}{x^2 - 3x + 2}$，则 $f^{(n)}(0) =$ \_\_\_\_.\par
% p141-b02; source PDF page 141
\blockanchor{p141-b02}
\question{4}不定积分 $\int |x|\mathrm{d}x =$ \_\_\_\_.\par
% p141-b03; source PDF page 141
\blockanchor{p141-b03}
\question{5}$\int_{0}^{+\infty}\frac{1}{(1 + x)^{2}(1 + x^{2024})}\mathrm{d}x =$\par
% p141-b04; source PDF page 141
\blockanchor{p141-b04}
\question{6}幂级数  $\sum_{n=0}^{\infty}(n+1)^{2}x^{n}$  的和函数为 \_\_\_\_.\par
% p141-b05; source PDF page 141
\blockanchor{p141-b05}
\question{7}若 $xyz e^{x + y + z} = 1$，则 $\frac{\partial^2z}{\partial x\partial y} =$ \_\_\_\_.\par
% p141-b06; source PDF page 141
\blockanchor{p141-b06}
\question{8}若广义积分 $\int_0^{+\infty} \frac{\sin(x^2)}{x^p} \, \mathrm{d}x$ 绝对收敛, 则 $p$ 的取值范围为 \_\_\_\_.\par
% p141-b07; source PDF page 141
\blockanchor{p141-b07}
\question{9}设 $[x]$ 表示不超过 $x$ 的最大整数, 则二重积分 $\iint_{D} [x + y] \mathrm{d} x \mathrm{d} y = \_\_\_\_$, 其中 $D = \{(x, y) \mid 0 \leq x \leq 3, 0 \leq y \leq 3\}$.\par
% p141-b08; source PDF page 141
\blockanchor{p141-b08}
\question{10}曲面积分 $\iint_{S} (x + y + z) \, \mathrm{d}S =$ \_\_\_\_，其中 $S$ 是球面 $x^{2} + y^{2} + z^{2} = a^{2}, z \geq 0$.\par
% p141-b09; source PDF page 141
\par\Needspace{5\baselineskip}
\blockanchor{p141-b09}
\majorquestion{二}计算题. 每题 15 分, 共 60 分.\par
% p141-b10; source PDF page 141
\par\Needspace{6\baselineskip}
\blockanchor{p141-b10}
\question{1}函数 $f(x)$ 在 $x = 1$ 附近有定义, 且在 $x = 1$ 处可导, 已知 $f(1) = 0$, $f'(1) = 2$. 求极限\par
\nopagebreak[4]
% p141-b11; source PDF page 141
\blockanchor{p141-b11}
\displayformula{\lim _ {x \rightarrow 0} \frac {f (\sin^ {2} x + \cos x)}{e ^ {x ^ {2}} - 1}.}
% p141-b12; source PDF page 141
\blockanchor{p141-b12}
\question{2}设 $D = \{(x, y) \mid x^2 + y^2 \leq 1\}$, 计算二重积分 $\iint_{D} \frac{\mathrm{d}x \mathrm{~d}y}{\sqrt{(1 - x)^2 + y^2}}$.\par
% p141-b13; source PDF page 141
\par\Needspace{6\baselineskip}
\blockanchor{p141-b13}
\question{3}求曲线积分\par
\nopagebreak[4]
% p141-b14; source PDF page 141
\blockanchor{p141-b14}
\displayformula{I = \int_ {L} (e ^ {x} + 1) \cos y \mathrm{d} x - [ (e ^ {x} + x) \sin y - x ] \mathrm{d} y.}
% p141-b15; source PDF page 141
\blockanchor{p141-b15}
其中 L 为由点 A(2,0) 沿着曲线  $y=\sqrt{4-x^{2}}$  到点 B(-2,0) 的有向曲线段.\par
% p141-b16; source PDF page 141
\blockanchor{p141-b16}
\question{4}将函数 $f(x) = \arctan \frac{1 - x}{1 + x}$ 展开成 $x$ 的幂级数, 并求级数 $\sum_{n=0}^{\infty} \frac{(-1)^n}{2n+1}$ 的和.\par
% p141-b17; source PDF page 141
\par\Needspace{5\baselineskip}
\blockanchor{p141-b17}
\majorquestion{三}证明题. 每题10分, 共30分.\par
% p141-b18; source PDF page 141
\par\Needspace{6\baselineskip}
\blockanchor{p141-b18}
\question{1}若函数 $f(x)$ 在 $[0,1]$ 上连续, $\int_0^1 f(x)\mathrm{d}x = 0$, 证明:\par
\nopagebreak[4]
% p141-b19; source PDF page 141
\blockanchor{p141-b19}
\displayformula{\left[ \int_ {0} ^ {1} x f (x) \mathrm{d} x \right] ^ {2} \leq \frac {1}{12} \int_ {0} ^ {1} f ^ {2} (x) \mathrm{d} x.}
% p141-b20; source PDF page 141
\blockanchor{p141-b20}
\question{2}证明函数 $f(x) = \sqrt{x} \ln x$ 在 $(0, +\infty)$ 上一致连续.\par
% p141-b21; source PDF page 141
\blockanchor{p141-b21}
\question{3}设函数 $f(x)$ 在 $[a, b]$ 上连续, 且 $f(x) > 0$. 记 $F(x) = \int_{a}^{b} f(t) |x - t| \, \mathrm{d}t$, 证明 $F(x)$ 在 $[a, b]$ 上有唯一的极小值点.\par
